\srcspaced{Jeder Zerfällungskörper von $\mA$ vom Grade $mr$ -- falls ein solcher existiert -- ist einem größten kommutativen Unterkörper von $\mA_r$ isomorph. Umgekehrt sind alle größten kommutativen Unterkörper von $\mA_r$ Zerfällungskörper, jeweils von einem Grad $ms$, wo $s$ ein Teiler von $r$ ist. Im regulären Fall sind alle größten kommutativen Unterkörper von $\mA_r$ vom Grade $mr$. Dabei wird unter regulärem Fall verstanden: Der Grundkörper $P$ hat die Eigenschaft, daß es zu jedem kommutativen Körper $\Sigma$ über $P$, der endlich über $P$, noch kommutative Erweiterungen über $\Sigma$ von beliebigem Grad gibt.}\footnote{Es liegt also beispielsweise kein regulärer Fall vor, wenn $P$ der Körper aller reellen Zahlen ist.}
